\section{Revisión del estado del arte y vacío de conocimiento}
\label{sec:estado-arte}

Esta revisión organiza la literatura en tres cuerpos que convergen sobre el
problema: la conceptualización de los gemelos digitales y los sistemas
ciberfísicos, su adopción específica en la industria de minerales, y la teoría
de estimación óptima de estados con sus aplicaciones al procesamiento de
minerales. El Anexo~\ref{anexo:matriz} presenta la matriz de revisión
bibliográfica con la caracterización sistemática de cada fuente.

\subsection{Gemelos digitales y sistemas ciberfísicos: la premisa implícita}

El paradigma de los sistemas ciberfísicos fue formulado por \textcite{lee2008}
como la integración estrecha entre procesos computacionales y procesos físicos,
señalando desde el inicio que el desafío central no es computacional sino de
fidelidad del acoplamiento entre ambos dominios. Sobre esa base,
\textcite{uhlemann2017} propusieron el gemelo digital como realización concreta
del sistema ciberfísico de producción para la Industria 4.0, y
\textcite{tao2019} lo orientaron hacia el pronóstico y la gestión de salud de
equipos complejos. \textcite{fuller2020} sistematizaron las tecnologías
habilitantes y los desafíos abiertos del área, y \textcite{liu2021} consolidaron
una revisión de conceptos, tecnologías y aplicaciones industriales.

Al comparar críticamente estos trabajos emerge una convergencia y una omisión
compartida. La convergencia es que todos coinciden en que la promesa del gemelo
digital reside en habilitar el mantenimiento predictivo y la toma de decisiones
autónoma. La omisión es que todos asumen, de manera tácita, la existencia de una
red de telemetría suficientemente limpia. \textcite{uhlemann2017} plantean la
integración de \emph{middleware} como vía para resolver las brechas de datos, es
decir, sitúan el problema en la capa de interoperabilidad y transporte;
\textcite{tao2019} y \textcite{fuller2020} centran su análisis en la
arquitectura y en la sincronización, no en la fidelidad de la señal de entrada.

\textcite{sepasgozar2021} es quien nombra con mayor precisión esta tensión al
distinguir el \emph{digital shadow} —una réplica que recibe datos crudos y los
espeja— del \emph{digital twin} propiamente tal, que procesa y depura la
incertidumbre subyacente. Esa distinción es la que este proyecto adopta como
criterio: sin una capa de depuración matemática explícita, lo que se construye
no es un gemelo digital sino una sombra digital de alta resolución.

\subsection{Gemelos digitales en la industria de minerales}

La traslación de este paradigma al dominio minero está documentada de manera
sistemática. \textcite{qu2023} desarrollan una revisión comprehensiva de los
gemelos digitales en la industria de minerales y evidencian una adopción
creciente pero fragmentada, con implementaciones concentradas en la
planificación y el transporte antes que en las etapas de conminución.

La revisión más reciente de \textcite{ebad2026}, que examina sesenta y ocho
estudios publicados entre 2015 y 2025 sobre gemelos digitales impulsados por
inteligencia artificial en minería, refuerza y actualiza ese diagnóstico. Los
autores reportan un predominio marcado de metodologías de aprendizaje automático
y aprendizaje profundo, e identifican como desafíos persistentes la integración
de datos, el procesamiento en tiempo real y —de manera notable para este
estudio— la escasa validación industrial de las propuestas.

La lectura conjunta de ambas revisiones permite una constatación relevante: el
esfuerzo investigativo en gemelos digitales mineros se ha concentrado en las
capas superiores de aprendizaje y predicción, mientras que la capa intermedia de
consistencia física de los datos permanece comparativamente desatendida. Se
avanza en la sofisticación del modelo predictivo sin avanzar proporcionalmente
en la garantía de calidad de sus entradas.

\subsection{Estimación óptima de estados y su aplicación en procesamiento de
minerales}

El marco formal para resolver la observabilidad parcial proviene de la teoría
clásica de estimación óptima. \textcite{kalman1960} estableció el estimador
recursivo de varianza mínima para sistemas lineales inmersos en ruido gaussiano
blanco, y \textcite{simon2006} consolidó el tratamiento moderno del problema,
incluyendo sus extensiones no lineales y las condiciones de convergencia
asintótica de la covarianza del error.

La aplicación de esta teoría al procesamiento de minerales tiene precedentes
sólidos y específicos. \textcite{apelt2002} desarrollaron la medición inferencial
de parámetros de molinos semiautógenos mediante estimación de estados,
demostrando que variables internas no instrumentadas del proceso de conminución
pueden reconstruirse a partir de mediciones indirectas. Más recientemente,
\textcite{barzaga2025} aplicaron filtros de Kalman extendido y \emph{unscented}
a la estimación de ley de cobre, consumo de ácido y contenido de humedad en
lixiviación en pilas, trabajo de particular pertinencia por su foco en la
minería del cobre y por provenir de un contexto operacional equivalente al
chileno.

El contraste entre estos dos grupos de literatura es el que estructura la
posición de este proyecto. La teoría de estimación óptima ofrece un instrumento
maduro y validado; la literatura de gemelos digitales mineros identifica una
carencia en la calidad y consistencia de los datos. Sin embargo, ambos cuerpos
se han desarrollado con escaso contacto entre sí: los trabajos de estimación de
estados en minería se formulan como problemas de control de procesos, y los
trabajos de gemelos digitales tratan la calidad del dato como un supuesto previo
antes que como un componente arquitectónico a diseñar.

\subsection{Identificación del vacío de conocimiento}

De la revisión anterior se desprende un vacío de conocimiento delimitado en
tres condiciones simultáneas.

\textbf{Primero}, existe evidencia de que la estimación de estados es viable y
efectiva en procesos de conminución \parencite{apelt2002} y en operaciones
mineras del cobre \parencite{barzaga2025}, pero ninguno de esos trabajos se
formula como capa de consistencia de un gemelo digital: se plantean como
instrumentos de control o de medición inferencial, con objetivos operacionales
distintos.

\textbf{Segundo}, la literatura de gemelos digitales industriales y mineros
reconoce la degradación e integración de datos como problema abierto
\parencite{fuller2020,ebad2026}, pero no incorpora el estimador estocástico como
elemento arquitectónico explícito entre la capa física y la capa de decisión.

\textbf{Tercero}, y de manera específica, no se identificó en la literatura
revisada un módulo analítico desacoplado, computacionalmente eficiente y
validado estadísticamente, que combine el principio físico de conservación de
masa con el Filtro de Kalman para reconstruir el volumen interno de carga de una
tolva de chancado primario bajo condiciones de vibración severa.

El vacío no consiste, por tanto, en la ausencia del algoritmo —que es teoría
establecida hace más de sesenta años— sino en la ausencia de su formulación como
componente arquitectónico de software, con alcance delimitado, validación
estadística formal y evidencia de reproducibilidad, dentro de una arquitectura
de gemelo digital minero. Precisar así el vacío es lo que permite que el aporte
del proyecto sea evaluable: no se reclama novedad algorítmica, sino aporte de
diseño, integración y validación.

\subsection{Arquitectura funcional del módulo propuesto}

La figura conceptual del flujo de información permite situar con exactitud el
componente que este proyecto desarrolla, destacado entre corchetes dobles:

\begin{center}
\singlespacing
\small
{}[Chancador físico] $\rightarrow$ \emph{(vibración y polvo)} $\rightarrow$
[Sensores / SCADA]\\
$\downarrow$\\
\textbf{Datos crudos degradados} $(y_k)$\\
$\downarrow$\\
\textbf{[[ Módulo de software analítico --- Filtro de Kalman (Python / NumPy) ]]}\\
$\downarrow$\\
\textbf{Estado estimado óptimo} $(\hat{\mathbf{x}}_{k|k}$ con MSE mínimo$)$\\
$\downarrow$\\{}
[Capa superior del gemelo digital / Simulación / Toma de decisiones]
\end{center}

Esta arquitectura funcional evidencia que los elementos seleccionados no son
piezas aisladas. El algoritmo de Kalman actúa como el motor matemático encargado
de depurar las señales corruptas del SCADA, garantizando que la réplica virtual
se construya sobre datos consistentes. El alcance de esta investigación se
concentra exclusivamente en el desarrollo e interacción del componente
destacado.
